% ============================================
% Chapter 7 – Conclusion and Future Work
% ============================================

\chapter{Conclusion and Future Work}

% -------------------------------------------------
% 7.1 Conclusion
% -------------------------------------------------

\section{Conclusion}

Building the Quantum Virtual Machine was an exercise in bridging the divide between theoretical quantum mechanics and practical software engineering. On paper, quantum algorithms operate on clean mathematical abstractions: qubits exist in pure superposition, gates perform instantaneous unitary rotations, and entanglement connects any pair of qubits without friction. In physical hardware, however, reality is messy: chips have fixed planar layouts, couplers are constrained, gates accumulate decoherence errors, and each hardware vendor promotes an incompatible software stack.

Our goal in this Final Year Project was to deliver a practical ``Write Once, Run Anywhere'' (WORA) runtime for quantum computing. Over the course of design, implementation, and benchmarking, we realized several concrete engineering milestones:

\begin{itemize}[leftmargin=*]
    \item \textbf{A Full Ingestion and Compilation Pipeline:} We developed an AST-based parser using Lark's formal grammar engine for OpenQASM 3.0, backed by parsers for OpenQASM 2.0 and JSON. The parser cleanly separates variable declarations from operation streams, handling classical memory registers, feed-forward conditions, and loops.
    \item \textbf{Topology-Aware Transpilation:} We proved that automated compiler heuristics can relieve developers from manual routing. By evaluating upcoming gates across an extended lookahead window with exponential decay weighting, our SABRE routing implementation cuts total SWAP overhead by 35\% to 50\% compared to standard greedy shortest-path routing, directly reducing circuit depth and noise vulnerability on target topologies.
    \item \textbf{Dual Simulation Backends:} Rather than relying on a single simulator, we engineered two complementary backends. The dense statevector engine uses multi-dimensional NumPy array slicing and axis permutations to execute single-qubit gates in microseconds without building expensive full-matrix Kronecker products. For larger systems, the Matrix Product State (MPS) simulator models states as 1D tensor trains, using Singular Value Decomposition (SVD) truncation to simulate 20+ qubits on consumer laptops as long as entanglement remains localized.
    \item \textbf{Realistic Noise Modeling:} To study how physical device imperfections degrade quantum logic, we built a Kraus operator open-system noise module. It models depolarizing white noise, $T_1$ thermal relaxation, $T_2$ pure dephasing, and asymmetric measurement readout errors, including pre-calibrated profiles of 5-qubit and 7-qubit physical devices.
    \item \textbf{Accessible Interfaces and Persistence:} We packaged QVM across three interfaces: a terminal CLI for scripting, an asynchronous FastAPI service with Supabase cloud database synchronization, and an interactive Next.js 16 / React 19 web dashboard that renders live SVG circuit diagrams, probability distributions, and KaTeX mathematical amplitudes.
    \item \textbf{Rigorous Verification:} The system was validated against 215 automated tests across 24 test suites, 26 high-depth rotation stress tests, and a 20-algorithm cross-framework audit that confirmed exact statevector equivalence with Qiskit and Cirq.
\end{itemize}

Ultimately, QVM demonstrates that hardware abstraction in quantum computing is both achievable and essential. By hiding the physical complexities of gate synthesis and coupling graphs behind a standardized runtime, developers and students can focus on the algorithmic logic that makes quantum computation powerful in the first place.

% -------------------------------------------------
% 7.2 Future Work
% -------------------------------------------------

\section{Future Work}

While QVM delivers a reliable testing ground for educational exploration and local algorithm verification, several technical directions warrant future investigation:

\textbf{Direct Cloud QPU Submission:} While QVM exports valid OpenQASM files compatible with third-party tools, adding native cloud submission adapters (e.g., using Qiskit Runtime or AWS Braket APIs) would let users submit transpiled circuits directly to physical dilution refrigerators from the QVM CLI or web dashboard.

\textbf{GPU-Accelerated Statevector Contraction:} The statevector engine currently runs on CPU using NumPy. Replacing NumPy tensor permutations with GPU-accelerated array libraries like CuPy or JAX would exploit parallel CUDA cores, potentially extending exact statevector simulation to 20--25 qubits on standard desktop workstations.

\textbf{Active Quantum Error Correction (QEC):} The current noise engine simulates open-system decay, but does not implement active syndrome measurement and error correction cycles. Integrating active simulations of the Surface code, Steane 7-qubit code, or Shor 9-qubit code would enable studying fault-tolerant threshold behavior under realistic noise.

\textbf{Higher-Dimensional Tensor Network Geometries:} The MPS engine assumes a one-dimensional chain of tensors, which fits linear ion traps or simple 1D topologies well. Extending the tensor network engine to Projected Entangled Pair States (PEPS) or tree tensor networks would allow simulating 2D grid chips with higher entanglement accuracy.

\textbf{Algorithmic Quantum Error Mitigation (QEM):} Adding post-processing error mitigation techniques, such as Zero-Noise Extrapolation (ZNE) and Clifford-based randomized benchmarking, would allow users to study how classical signal processing can extract cleaner expectation values from noisy simulation runs.

% -------------------------------------------------
% 7.3 Limitations and Known Weaknesses
% -------------------------------------------------

\section{Limitations and Known Weaknesses}

To assess the Quantum Virtual Machine objectively, one must evaluate the mathematical boundaries of classical quantum simulation alongside the specific architectural trade-offs made during its implementation. Acknowledging these constraints provides a realistic appraisal of QVM's capabilities and establishes clear technical benchmarks for future engineering.

\subsection{Computational and Algorithmic Weaknesses}

\begin{itemize}[leftmargin=*]
    \item \textbf{Classical Memory Scaling and Statevector RAM Ceiling:} Classical computers cannot cheat the exponential scaling of quantum mechanics. A dense statevector requires $2^n$ complex numbers (16 bytes each in double-precision \texttt{complex128}). While 10 qubits consumes only 16~KB, memory demands escalate steeply: 14 qubits requires 256~KB, 16 qubits requires 1~MB, 20 qubits requires 16~MB, and 30 qubits demands 16~GB of contiguous memory. In Python and NumPy, tensor reshaping and axis permutations create transient memory allocations up to three times the statevector size during multi-qubit gate operations. Consequently, unassisted desktop execution encounters a practical ceiling around 14 qubits on typical consumer hardware before operating system paging degrades performance.
    \item \textbf{Volume-Law Entanglement Breakdown in Matrix Product States:} While the MPS engine simulates weakly entangled circuits up to 20+ qubits by decomposing states into 1D tensor trains, its tractability depends strictly on the decay rate of Schmidt singular values across bipartite cuts. In circuits that generate global, volume-law entanglement---such as the Quantum Fourier Transform (QFT), deep random circuit sampling, or all-to-all variational ansatzes---the required bond dimension $\chi$ scales exponentially ($\chi \sim 2^{n/2}$). Under fixed SVD truncation limits ($\chi_{\text{max}} = 64$), the engine discards non-negligible entanglement amplitudes, causing state fidelity $\mathcal{F} = |\langle \psi_{\text{exact}} | \psi_{\text{MPS}} \rangle|^2$ to degrade below acceptable thresholds.
    \item \textbf{Sub-optimality of Heuristic Lookahead Routing:} Mapping logical qubits onto arbitrary hardware coupling graphs is an NP-hard combinatorial search problem. While QVM's SABRE implementation reduces total SWAP gate counts by 35\% to 50\% relative to naive greedy BFS, it remains a greedy lookahead heuristic operating across an extended gate window $\mathcal{E}$. On irregular topologies, cyclic coupling graphs, or processors with non-uniform two-qubit gate fidelities, the heuristic cost function can become trapped in local minima, introducing redundant SWAP cascades that unnecessarily inflate circuit depth.
    \item \textbf{Stochastic Sampling Latency in Kraus Noise Trajectories:} Full density matrix evolution scales as $\mathcal{O}(4^n)$, which exhausts classical memory even faster than statevectors ($n=10$ requires $10^6$ complex entries). QVM resolves this by employing Monte Carlo quantum trajectory sampling. However, this introduces a runtime trade-off: simulating $N_{\text{shots}} = 10{,}000$ requires $10{,}000$ independent stochastic passes through the circuit. Under high gate depths, this produces linear execution time dilation ($T \propto N_{\text{shots}} \times \text{depth}$) compared to deterministic unitary statevector evolution.
\end{itemize}

\subsection{Architectural and Practical Delimitations}

\begin{itemize}[leftmargin=*]
    \item \textbf{Absence of Direct Cloud QPU Dispatch Brokers:} QVM compiles circuits into syntactically valid OpenQASM 2.0 and 3.0 representations that strictly conform to IBM Quantum and Rigetti hardware specifications. However, the virtual machine currently lacks built-in cloud broker credentials, OAuth token refreshes, and automated remote queue polling. Users must manually copy transpiled OpenQASM artifacts into third-party SDKs (such as \texttt{qiskit-ibm-runtime} or \texttt{amazon-braket-sdk}) to dispatch jobs to physical dilution refrigerators.
    \item \textbf{Lack of Active Stabilizer Fault-Tolerant Decoding:} QVM models physical noise phenomenologically via Kraus operator channels (depolarizing white noise, $T_1$ thermal relaxation, $T_2$ pure dephasing, and readout assignment error). It does not implement real-time stabilizer syndrome extraction loops or surface-code minimum-weight perfect matching (MWPM) decoders. The platform simulates the destructive effects of noise on physical circuits, but cannot actively execute fault-tolerant error correction cycles.
    \item \textbf{Single-Node CPU Execution Boundaries:} The core simulation engines rely on single-node vectorized NumPy execution without distributed memory parallelization (such as MPI across multiple compute nodes) or GPU acceleration (CUDA kernels). This bounds the maximum addressable Hilbert space to the physical memory of a single host machine.
\end{itemize}

Despite these physical and classical computational boundaries, QVM achieves its primary educational and architectural mission: providing a transparent, lightweight, and hardware-independent runtime for quantum computing.
